\chapter{Related Work}\label{ch:related}

This chapter positions the thesis within prior work on split neural-network inference,
cloud-native deployment overhead, and security hardening for microservice systems. The
central distinction is that most prior partitioning work focuses on selecting or
optimising split points, while this thesis follows one workload through a staged
software-engineering evaluation: local split screening, Kubernetes deployment, managed
AKS transfer validation, and security hardening.

\section{Split Computing and DNN Partitioning}
\label{sec:rw-split-computing}

Early split-computing systems established that neural-network inference can be divided
between execution environments, but that the split point determines whether the
distributed design is useful. Neurosurgeon models the latency trade-off between mobile
and cloud execution and shows that different layers have different computation and data
transfer profiles~\cite{neurosurgeon}. JointDNN frames partitioning as an optimisation
problem in which upload and download costs depend on the intermediate tensor at the
chosen boundary~\cite{jointdnn}. DeeperThings extends this line of work to fully
distributed CNN inference on constrained edge devices and reinforces the observation
that early feature maps can be expensive to move because they remain spatially
large~\cite{deeperthings}.

More recent work broadens the partitioning objective beyond a single latency value.
DNNSplit selects split points by jointly considering latency and cost in multi-tier
settings~\cite{dnnsplit}. Genetic-algorithm and delay-aware approaches likewise treat
the split point as dependent on device capability, bandwidth, throughput, and parallelism
constraints~\cite{ga_partitioning,delay_aware_inference}. Survey work on DNN
partitioning and split computing summarizes this literature as a multi-objective problem
involving activation size, compute allocation, deployment tier, and runtime
conditions~\cite{survey_partitioning}. This thesis adopts that view but keeps the
optimisation scope deliberately narrow: it studies a fixed ResNet-18 workload and
evaluates how a small number of architecturally meaningful split choices behave across
deployment layers.

\section{Granularity and Internal Model Structure}
\label{sec:rw-granularity}

Several studies show that coarse layer depth is not enough to explain split performance.
Fine-grained elastic partitioning demonstrates that more detailed placement decisions
can improve distributed DNN inference when the placement is aware of both computation
and communication conditions~\cite{fine_grained_partitioning}. Hierarchical
partitioning similarly combines global partitioning with local refinement, showing
that staged decisions can account for heterogeneity at both cluster and node
levels~\cite{hierarchical_partitioning}.
Pareto-front analysis of DNN partitioning further shows that execution time can vary
across blocks inside the same model, so boundaries with similar activation sizes may
still differ in compute distribution~\cite{pareto_split}. BottleFit and BottleNet++
approach the problem from the feature-compression side, showing that intermediate
representations can be compressed or shaped to reduce communication cost, but that the
position of the boundary remains central to the latency and accuracy
trade-off~\cite{bottlefit,bottlenet_pp}.

The experiments in this thesis use these findings to justify a coarse-to-fine design.
Stage~L1 evaluates residual-stage boundaries in ResNet-18, and Stage~L2 adds a block-level
boundary inside the carried-forward region. The goal is not to claim that the selected
boundaries are globally optimal for all DNNs. Instead, the goal is to show how
activation-transfer size and compute placement interact in a reproducible microservice
setting.

\section{Kubernetes and Cloud-Native Inference}
\label{sec:rw-cloud-native}

Moving from local split inference to Kubernetes changes the systems context. A
microservice boundary is not only a model partition; it is also an RPC path with
serialization, protocol handling, socket operations, scheduling effects, and service
networking. Recent RPC and microservice studies show that communication overhead can
come from layered protocol processing and network-stack traversal, not only from the
application payload itself~\cite{rpc_overhead_hotos25,notnets_apsys24}. Service-mesh
work reaches a similar conclusion for proxied microservices: sidecars and data-plane
components add workload-dependent overhead through additional critical-path
processing~\cite{meshinsight_socc23}.

Kubernetes and cloud studies also show why local and managed-cluster results should be
compared carefully. Managed Kubernetes performance varies across services and
configurations~\cite{managed_k8s_perf}. Container and VM layers can differ from
bare-metal execution~\cite{virtualization_costs}. CNI plugins and edge-oriented
Kubernetes networking can affect communication cost~\cite{cni_k8s_ic2e21}. Cloud
environments also exhibit network-performance variability that can change application
scalability and measurement stability~\cite{noise_in_clouds}. For this reason, this
thesis does not treat the local Kubernetes and AKS stages as numerical replicas.
Instead, Stage~C asks whether the ordering and relative overhead trends from local
Kubernetes transfer directionally to managed AKS.

\section{Security Hardening}
\label{sec:rw-security}

The security stages build on prior work in service identity, mTLS, service meshes, and
confidential computing. Policy-driven Kubernetes security work motivates combining
encrypted communication with explicit authorization policy rather than treating TLS
alone as the complete protection mechanism~\cite{ashraf2025policy_workflow}. Service
mesh studies show that sidecar-based mTLS can improve service-to-service trust
boundaries, but also that it can add latency and CPU/resource overhead, with the
cost depending on mesh configuration and workload~\cite{meshinsight_socc23,bremlerbarr2025mtls}.
Work on service-mesh control-plane
trust further cautions that mesh CAs and control planes remain security-critical
components~\cite{poudel2025mazu}. Stage~H1 therefore treats mTLS and authorization as a
bounded communication-hardening layer, not as a complete security solution.

Confidential-computing work motivates the second hardening layer. The motivation
itself comes from work on activation-level privacy attacks:
He~et~al.~\cite{he2021_attacking_collaborative} demonstrate that a remote partition
holding only intermediate feature maps can reconstruct the edge client's input, and
that simple noise-based defenses fail. This shows that in-transit mTLS alone does
not protect activations once they reach the downstream segment, although their
threat model is malicious-remote-service rather than the infrastructure-level
adversary that confidential VMs address.

AMD SEV-SNP provides VM-level memory encryption and integrity protection against
parts of the host platform~\cite{amd_sevsnp_whitepaper,kaplan2023_hardware_vm}.
Comparative studies of virtualization-based confidential-computing platforms
emphasize that these systems have specific attacker models, attestation
mechanisms, and residual limitations~\cite{guanciale2022_confidential_quartet,misono2024_cvm_explained}.
Performance studies show that confidential VM overhead is workload-dependent,
with CPU-bound, memory-intensive, and I/O-heavy workloads affected
differently~\cite{misono2024_cvm_explained,akram2021_hpc_tees}. Machine-learning
confidential computing surveys identify protected model execution and data
privacy as important motivations, while also emphasizing deployment and
performance trade-offs~\cite{mo2024_ml_confidential_computing}. Recent work
additionally shows that the hypervisor remains a meaningful adversary even
under SEV-SNP: HECKLER~\cite{schluter2024_heckler} uses malicious-interrupt
injection to bypass OpenSSH and sudo authentication on production SEV-SNP
and TDX, and additionally breaks execution integrity of analytical
workloads on SEV-SNP, an attack surface that hardware memory encryption
alone does not cover. Stage~H2 follows this bounded framing by
measuring VM-level confidential execution on AKS for the downstream service
only, rather than claiming application-level enclave isolation or evaluating
those residual attack surfaces.

\section{Research Gap Analysis and Scope}
\label{sec:rw-positioning}

The related literature establishes the individual concerns that motivate this thesis:
split points affect activation-transfer and compute costs, microservice and Kubernetes
deployment add platform overhead, and security hardening strengthens protection while
introducing additional cost. Each concern, however, is typically studied in isolation.
Partitioning work optimises split points without carrying the chosen decomposition into
an orchestrated or hardened deployment, while cloud-native and service-mesh work
characterises platform and security overhead without tying it to a specific
model-partitioning decision. The gap this thesis addresses is the absence of a single
staged evaluation that follows one workload across all of these layers. Instead of
evaluating partitioning, orchestration, and hardening as separate problems, the thesis
tracks the same ResNet-18 microservice inference workload from local split selection
through Kubernetes, AKS, mTLS/AuthZ, and confidential VM execution.

To make the link between the surveyed gaps and the research questions explicit, the
following summarises which gap each research question --- and its constituent stages ---
addresses:

\begin{description}
    \item[\textbf{RQ1} (architectural split \& deployment).] Prior partitioning work
    selects or optimises split points, frequently for a single execution environment
    (\cref{sec:rw-split-computing,sec:rw-granularity}), while cloud-native studies
    characterise deployment overhead independently of any particular partition
    (\cref{sec:rw-cloud-native}). RQ1 connects these by screening coarse and fine
    boundaries on a fixed ResNet-18 workload under controlled local execution
    (Stages~L1--L3) and then carrying the \emph{same} decomposition into local Kubernetes
    (Stage~K) and managed AKS (Stage~C), testing whether the split behaviour established
    locally transfers directionally to orchestrated and cloud settings rather than
    assuming that it does.
    \item[\textbf{RQ2} (security hardening).] Service-mesh and confidential-computing work
    shows that mutual TLS and VM-level protection add workload-dependent overhead
    (\cref{sec:rw-security}), but typically in isolation from a concrete partitioned
    inference deployment. RQ2 connects these by measuring the incremental performance and
    operational cost of communication-level hardening (Stage~H1) and selective VM-level
    confidential execution (Stage~H2) on the \emph{specific} AKS deployment selected under
    RQ1, and by synthesising the result into a threat-model-aware trade-off recommendation.
\end{description}

Across both questions, the distinguishing scope of this thesis is therefore not a new
partitioning algorithm or a new security mechanism, but the end-to-end, single-workload
staged pipeline that connects split selection, deployment context, and security hardening
under one reproducible measurement contract.
